\section*{Lowpass VCVS Butterworth}

\noindent\colorbox{orange}{
  \begin{minipage}{1.0\linewidth}
    \paragraph{Richiesta (a)}
    Dimostrare, sulla base di semplici argomenti di natura qualitativa, che il circuito proposto
    svolga effettivamente la funzione di filtro passa-basso, discutendone la risposta ai limiti
    del dominio delle frequenze reali.
  \end{minipage}
}

\paragraph{Risposta (a)}
\begin{itemize}
\item Per $f \to 0$
  i condensatori agiscono da rami aperti, le resistenze hanno impedenze infinite che vengono
  assorbite dall'impedenze ideali all'ingresso dell'opamp. Questa \'e la configurazione di
  un voltage follower. La tensione d'uscita \'e la tensione d'ingresso, il guadagno \'e unitario.
  
\item Per $f \to \infty$
  i condensatori agiscono da cortocircuiti, l'ingresso non-invertente viene messo a terra,
  l'ingresso invertente nell'ipotesi di corto-circuito virtuale deve eguagliare la tensione,
  la tensione di uscita deve essere nulla.
\end{itemize}


\noindent\colorbox{orange}{
  \begin{minipage}{1.0\linewidth}
    \paragraph{Richiesta (b)}
    Derivare analiticamente la funzione di trasferimento complessa $A_v(s)=V_{OUT}(s)/V_S(s)$
    nell’ipotesi che l’operazionale sia ideale (suggerimento: si determini la tensione nel nodo
    A del circuito in termini di $V_S$ e $V_{OUT}$; quindi si esprima $V_A$ in termini di $V_{OUT}$ assumendo
    che l'opAmp funzioni in regime lineare ed utilizzando il principio di cortocircuito
    virtuale; si eguaglino infine le espressioni ottenute per ottenere la funzione di
    trasferimento). Per semplicit\'a di calcolo, si assuma inoltre (ipotesi compatibile con i
    valori e le incertezze nominali dei componenti) che R1=R2=R, C1=C, C2=2C.
  \end{minipage}
}

\paragraph{Risposta (b)}
Ipotesi: l'operazionale \'e ideale ed opera in regime lineare.
\begin{itemize}
\item
  Tra $V_-$ e $V_{OUT}$ i rami sono corto-circuiti, di conseguenza $V_- = V_{OUT}$.
\item
  Applico il principio di \textit{cortocircuito virtuale} e di conseguenza $V_+ = V_-$.
\end{itemize}
di conseguenza
\begin{equation*}
  V_B = V_{OUT}
\end{equation*}
le successive equazioni nodali appaiate con il risultato precedente permettono di
ottenere la funzione di trasferimento corretta
\begin{equation*}
  \begin{gathered}
    \frac{V_A - V_B}{R} = sCV_{OUT} \iff V_A = (1 + sRC)V_{OUT} \\
    \frac{V_S - V_A}{R} = \frac{V_A - V_B}{R} + 2sC(V_A - V_{OUT}) \\
    V_S - V_A = V_A - V_B + 2sRC(V_A - V_{OUT}) \\
    V_S = (2 + 2sRC)V_A - (1 + 2sRC)V_{OUT} \\
    V_S = \left[(2 + 2sRC)(1 + sRC) - (1 + 2sRC)\right]V_{OUT} \\
  \end{gathered}
\end{equation*}
da cui
\begin{equation}
  \label{eq:lp-vcvs-butterworth-tfunc}
  \colorbox{blu}{\boxed{H(s) = \frac{1}{1 + 2RCs + 2R^2C^2s^2}}}
\end{equation}

\noindent\colorbox{orange}{
  \begin{minipage}{1.0\linewidth}
    \paragraph{Richiesta (c)}
    Dalla funzione di trasferimento dedurre le caratteristiche del filtro osservate al punto 1)
    (guadagno a centro-banda, frequenza di taglio, andamento di modulo/fase in funzione
    della frequenza) e confrontarle con quanto misurato ai punti 1c, 1d, 1e.
  \end{minipage}
}

\paragraph{Risposta (c)}
il guadagno \'e ottenuto per mezzo dell'esspressione \ref{eq:da-tfunc-a-gain}
\begin{equation*}
  \begin{gathered}
    \begin{cases}
      \mathcal{R}_{\omega}^N = 1 & \mathcal{I}_{\omega}^N = 0 \\
      \mathcal{R}_{\omega}^D = 1 - 2R^2C^2\omega^2 & \mathcal{I}_{\omega}^D = -2RC\omega
    \end{cases} \\
    G(\omega) = \frac{1}{\sqrt{(1 - 2R^2C^2\omega^2)^2 + (-2RC\omega)^2}} = \frac{1}{\sqrt{1 + (4R^4C^4)\omega^4}}
  \end{gathered}
\end{equation*}
\begin{equation}
  \label{eq:lp-vcvs-butterworth-gain}
  \colorbox{blu}{\boxed{G(f) = \frac{1}{\sqrt{1 + (64\pi^4R^4C^4)f^4}}}}
\end{equation}
per determinare la fase si usa \ref{eq:da-tfunc-a-fase-immaginario-negativo}
\begin{equation*}
  \phi(\omega) = \arctan\left(\frac{1 - 2R^2C^2\omega^2}{2RC\omega}\right)
\end{equation*}
\begin{equation}
  \label{eq:lp-vcvs-butterworth-phase}
  \colorbox{blu}{\boxed{\phi(f) = -\frac{\pi}{2} + \arctan\left(\frac{1 - 8\pi^2R^2C^2f^2}{4\pi RCf}\right)}}
\end{equation}


Frequenza di taglio.
\begin{equation*}
  \begin{gathered}
    G(f) = \frac{\max\{G(\omega)\}}{\sqrt2} = \frac{1}{\sqrt2} \implies \frac{1}{\sqrt{1 + (64\pi^4R^4C^4)f^4}} = \frac{1}{\sqrt2} \\
  \end{gathered}
\end{equation*}
\begin{equation}
  \label{eq:lp-vcvs-butterworth-fcut}
  \colorbox{blu}{\boxed{f = \frac{1}{2\pi\sqrt2\cdot RC}}}
\end{equation}

\paragraph{Richiesta (d)}
Dare un’interpretazione qualitativa (e se possibile quantitativa) delle caratteristiche del
segnale di uscita osservato al punto 1f. Che tipo di segnale appare? A quale ampiezza?
